% Generated draft: requires mathematical and source review.
\documentclass[11pt]{article}
\usepackage[margin=25mm]{geometry}
\usepackage{fontspec}
\setmainfont{Latin Modern Roman}
\usepackage{amsmath,amssymb,mathtools,unicode-math}
\setmathfont{Latin Modern Math}
\usepackage{enumitem,needspace,xurl,hyperref}
\setlength{\parindent}{0pt}
\setlength{\parskip}{5pt}
\setlength{\emergencystretch}{3em}
\newcommand{\N}{\mathbb{N}}
\newcommand{\Z}{\mathbb{Z}}
\newcommand{\Q}{\mathbb{Q}}
\newcommand{\R}{\mathbb{R}}
\newcommand{\C}{\mathbb{C}}
\newcommand{\F}{\mathbb{F}}
\newcommand{\A}{\mathbb{A}}
\newcommand{\PP}{\mathbb{P}}
\newcommand{\E}{\mathbb{E}}
\newcommand{\eps}{\varepsilon}
\newcommand{\dd}{\,\mathrm{d}}
\begin{document}
\section*{UnsolvedMath ID: 30004686}
\textbf{Source problem code:} catalogue entry
\subsection{Definitions and mathematical statement}
An evolution equation governs the time dependence of an unknown mathematical object. The supplied record does not specify the unknown, the class of equations, the spatial domains, or the initial and boundary conditions. It also does not define what dimensions count as higher-dimensional or what properties constitute a unifying theory.

The source poses the following research question:
\begin{quote}
Can there be a unifying theory for the analytical and numerical treatment of evolution equations on higher-dimensional domains?
\end{quote}
The background and supplied literature assessment restrict attention to \emph{dissipative} evolution equations, whereas the recorded statement does not state that restriction. Neither a definition of dissipativity nor a criterion for a successful analytical and numerical theory is supplied. Thus this is a research programme whose precise scope remains unresolved; the evidence does not determine a mathematical assertion with explicit hypotheses and quantifiers.
\subsection{Short English statement}
The question asks for a common framework for analysing evolution equations and treating them numerically on higher-dimensional domains. The background specifically concerns dissipative equations. The intended equation class, dimensions, assumptions, and required conclusions remain unspecified, so a precise theorem or conjecture cannot yet be stated.
\subsection{Sources}
\begin{itemize}
\item Local UnsolvedMath registry, \texttt{data/unsolvedmath/problems.json}: permanent ID \(30004686\), title \emph{Unified Analysis of Higher-Dimensional Evolution Equations}, source problem code \texttt{catalogue entry}; collection: Oberwolfach Reports Open Problems. No individual proposer is supplied; the proposed year is recorded as 2021.
\item Supplied report citation: \emph{Mini-Workshop: Mathematics of Dissipation -- Dynamics, Data and Control} (2022), Oberwolfach Reports, pp.\ 1259--1289, DOI \url{https://doi.org/10.4171/owr/2021/24}. The supplied literature-source label is \emph{Oberwolfach Report 24/2021}, with the same URL. The background describes these reports as records of mathematical workshops and their open-problem discussions. No specific page, passage, or named contributor for this question is supplied.
\item Literature assessment dated \texttt{2026-08-22}: the supplied curation labels the problem open and states that the report identifies a unifying analytic and numerical theory for dissipative evolution equations in higher-dimensional domains as an open problem, with existing results covering important classes but not the requested general theory. This is supplied dated curation, not independently verified. Its source-status and literature claims require independent verification.
\item The registry records statement status \texttt{corrected\_verified} and reports a manual semantic, typographic, and LaTeX review concluding that the clean statement is self-contained and preserves the target without unsupported additions. This is a supplied editorial assessment; it does not resolve the formulation gaps identified here or constitute independent mathematical review.
\end{itemize}
\end{document}
